\sectionkicker{Source-backed technical notes}
\subsection*{フロンティアモデルは「会話するモデル」からマルチモーダルな対話基盤へ: Technical Notes}
本欄は記事本文で圧縮した一次資料上の情報を、比較・再検証しやすい形で整理したものである。時系列、確認済みの主張、留意点、一次資料URLを掲載する。
\smallskip
\noindent{\small\textit{Technical Notesでは、一次資料から確認した公開・提供・研究上の事実を記載する。数値・能力評価や提供元・著者による比較表現は、独立再現済みの結果ではなく帰属claimとして扱う。}}\par
\medskip
\subsection*{Theme at a glance}
\begin{center}
\footnotesize
\begin{tabularx}{\linewidth}{@{}p{0.20\linewidth}p{0.10\linewidth}p{0.14\linewidth}X@{}}
\toprule
資料 & 位置づけ & 種別 & 時系列 \\
\midrule
GPT-4o & 主要資料 & モデル & 2024-05-13     \\
Gemini 1.5 / Gemini 1.5 Flash & 主要資料 & モデル更新 & 2024-02-15, 2024-05-14     \\
Claude 3 → Claude 3.5 Sonnet & 主要資料 & モデル更新 & 2024-06-21, 2024-03-04     \\
Quiet-STaR & 補足資料 & 論文 & 2024-03-14     \\
Infini-attention & 補足資料 & 論文 & 2024-04-10     \\
Chameleon & 補足資料 & モデル & 2024-05-16, 2024-06-18     \\
\bottomrule
\end{tabularx}
\end{center}
\normalsize

\begin{technicalnote}{GPT-4o}{主要資料}
\begingroup
% reader-facing Technical Notes break policy
\widowpenalty=10000
\clubpenalty=10000
\displaywidowpenalty=10000
\begin{tabularx}{\linewidth}{@{}>{\bfseries}p{0.22\linewidth}X@{}}
組織 & OpenAI \\
種別 & モデル \\
時系列 & 2024-05-13 \\
\end{tabularx}
\smallskip
\begin{minipage}{\linewidth}
% reader-facing Technical Notes event-only fact/source tail group
{\bfseries 一次資料から整理したtechnical points}
\begin{itemize}[leftmargin=1.5em,itemsep=0.35em]
\item \textbf{一次情報で確認できる事実}: OpenAIの一次資料により、GPT-4oについて2024-05-13にモデル公開を確認した。 GPT-4oはtext・vision・audioを単一neural networkでend-to-end学習し、従来Voice Modeのspeech-to-text→LLM→text-to-speechという3-model pipelineでは失われていたtone・multiple speakers・background noise等を直接扱う設計として発表された。2024年5月13日のlaunch時点で公開されたのは主にtext/image inputとtext outputで、audio/video capabilityは段階提供予定だったため、model architecture上のmultimodalityと当日のproduct availabilityを分けて扱う。benchmark・latency・price比較はOpenAIによる評価として扱う。
\end{itemize}
\begin{minipage}{\linewidth}
% reader-facing Technical Notes generic boundary/source tail group
\begin{samepage}
{\bfseries 一次資料}
\begingroup\sloppy
\Urlmuskip=0mu plus 2mu\relax
\begin{itemize}[leftmargin=1.5em,itemsep=0.25em]
\item {\scriptsize\url{https://openai.com/index/hello-gpt-4o}}
\end{itemize}
\endgroup
\end{samepage}
\end{minipage}
\end{minipage}
\endgroup
\end{technicalnote}
\medskip

\begin{technicalnote}{Gemini 1.5 / Gemini 1.5 Flash}{主要資料}
\begingroup
% reader-facing Technical Notes break policy
\widowpenalty=10000
\clubpenalty=10000
\displaywidowpenalty=10000
\begin{tabularx}{\linewidth}{@{}>{\bfseries}p{0.22\linewidth}X@{}}
組織 & Google \\
種別 & モデル更新 \\
時系列 & 2024-02-15, 2024-05-14 \\
\end{tabularx}
\smallskip
\begin{minipage}{\linewidth}
% reader-facing Technical Notes event-only fact/source tail group
{\bfseries 一次資料から整理したtechnical points}
\begin{itemize}[leftmargin=1.5em,itemsep=0.35em]
\item \textbf{一次情報で確認できる事実}: Googleの一次資料により、Gemini 1.5 / Gemini 1.5 Flashについて2024-02-15にモデルPreview、2024-05-14にモデル発表を確認した。 対象event近傍の一次資料から Mixture-of-Experts (MoE) を確認できる。
\end{itemize}
\begin{minipage}{\linewidth}
% reader-facing Technical Notes generic boundary/source tail group
\begin{samepage}
{\bfseries 一次資料}
\begingroup\sloppy
\Urlmuskip=0mu plus 2mu\relax
\begin{itemize}[leftmargin=1.5em,itemsep=0.25em]
\item {\scriptsize\url{https://blog.google/innovation-and-ai/products/google-gemini-next-generation-model-february-2024/}}
\item {\scriptsize\url{https://blog.google/innovation-and-ai/technology/developers-tools/gemini-gemma-developer-updates-may-2024/}}
\end{itemize}
\endgroup
\end{samepage}
\end{minipage}
\end{minipage}
\endgroup
\end{technicalnote}
\medskip

\begin{technicalnote}{Claude 3 → Claude 3.5 Sonnet}{主要資料}
\begingroup
% reader-facing Technical Notes break policy
\widowpenalty=10000
\clubpenalty=10000
\displaywidowpenalty=10000
\begin{tabularx}{\linewidth}{@{}>{\bfseries}p{0.22\linewidth}X@{}}
組織 & Anthropic \\
種別 & モデル更新 \\
時系列 & 2024-06-21, 2024-03-04 \\
\end{tabularx}
\smallskip
\begin{minipage}{\linewidth}
% reader-facing Technical Notes event-only fact/source tail group
{\bfseries 一次資料から整理したtechnical points}
\begin{itemize}[leftmargin=1.5em,itemsep=0.35em]
\item \textbf{一次情報で確認できる事実}: Anthropicの一次資料により、Claude 3 → Claude 3.5 Sonnetについて2024-03-04にモデル公開、2024-06-21にモデル公開を確認した。 対象event近傍の一次資料から 200K context を確認できる。
\end{itemize}
\begin{minipage}{\linewidth}
% reader-facing Technical Notes generic boundary/source tail group
\begin{samepage}
{\bfseries 一次資料}
\begingroup\sloppy
\Urlmuskip=0mu plus 2mu\relax
\begin{itemize}[leftmargin=1.5em,itemsep=0.25em]
\item {\scriptsize\url{https://www.anthropic.com/news/claude-3-5-sonnet}}
\item {\scriptsize\url{https://www.anthropic.com/news/claude-3-family}}
\end{itemize}
\endgroup
\end{samepage}
\end{minipage}
\end{minipage}
\endgroup
\end{technicalnote}
\medskip

\begin{technicalnote}{Quiet-STaR}{補足資料}
\begingroup
% reader-facing Technical Notes break policy
\widowpenalty=10000
\clubpenalty=10000
\displaywidowpenalty=10000
\begin{tabularx}{\linewidth}{@{}>{\bfseries}p{0.22\linewidth}X@{}}
組織 & - \\
種別 & 論文 \\
時系列 & 2024-03-14 \\
\end{tabularx}
\smallskip
\begin{minipage}{\linewidth}
% reader-facing Technical Notes event-only fact/source tail group
{\bfseries 一次資料から整理したtechnical points}
\begin{itemize}[leftmargin=1.5em,itemsep=0.35em]
\item \textbf{一次情報で確認できる事実}: 提供元の一次資料により、Quiet-STaRについて2024-03-14にPaper First Submissionを確認した。 Quiet-STaRは、将来のテキストを説明するrationaleを各token位置で生成する学習法として、tokenwise parallel sampling、thoughtの開始・終了を示すlearnable token、extended teacher forcingを組み合わせる。著者らはcontinued pretraining後のGSM8K・CommonsenseQA等の改善を報告しているが、数値性能は著者報告として扱う。
\end{itemize}
\begin{minipage}{\linewidth}
% reader-facing Technical Notes generic boundary/source tail group
\begin{samepage}
{\bfseries 一次資料}
\begingroup\sloppy
\Urlmuskip=0mu plus 2mu\relax
\begin{itemize}[leftmargin=1.5em,itemsep=0.25em]
\item {\scriptsize\url{http://arxiv.org/abs/2403.09629v2}}
\end{itemize}
\endgroup
\end{samepage}
\end{minipage}
\end{minipage}
\endgroup
\end{technicalnote}
\medskip

\begin{technicalnote}{Infini-attention}{補足資料}
\begingroup
% reader-facing Technical Notes break policy
\widowpenalty=10000
\clubpenalty=10000
\displaywidowpenalty=10000
\begin{tabularx}{\linewidth}{@{}>{\bfseries}p{0.22\linewidth}X@{}}
組織 & Google \\
種別 & 論文 \\
時系列 & 2024-04-10 \\
\end{tabularx}
\smallskip
\begin{minipage}{\linewidth}
% reader-facing Technical Notes event-only fact/source tail group
{\bfseries 一次資料から整理したtechnical points}
\begin{itemize}[leftmargin=1.5em,itemsep=0.35em]
\item \textbf{一次情報で確認できる事実}: Googleの一次資料により、Infini-attentionについて2024-04-10にPaper First Submissionを確認した。 Infini-attentionは、通常のattention blockへcompressive memoryを組み込み、masked local attentionとlong-term linear attentionを単一Transformer block内で扱う構成である。著者らはbounded memory/computationを保った長文脈処理とstreaming inferenceを目的として提示しており、長文脈タスク上の性能値は著者報告として扱う。
\end{itemize}
\begin{minipage}{\linewidth}
% reader-facing Technical Notes generic boundary/source tail group
\begin{samepage}
{\bfseries 一次資料}
\begingroup\sloppy
\Urlmuskip=0mu plus 2mu\relax
\begin{itemize}[leftmargin=1.5em,itemsep=0.25em]
\item {\scriptsize\url{http://arxiv.org/abs/2404.07143v2}}
\end{itemize}
\endgroup
\end{samepage}
\end{minipage}
\end{minipage}
\endgroup
\end{technicalnote}
\medskip

\begin{technicalnote}{Chameleon}{補足資料}
\begingroup
% reader-facing Technical Notes break policy
\widowpenalty=10000
\clubpenalty=10000
\displaywidowpenalty=10000
\begin{tabularx}{\linewidth}{@{}>{\bfseries}p{0.22\linewidth}X@{}}
組織 & Meta \\
種別 & モデル \\
時系列 & 2024-05-16, 2024-06-18 \\
\end{tabularx}
\smallskip
\begin{minipage}{\linewidth}
% reader-facing Technical Notes event-only fact/source tail group
{\bfseries 一次資料から整理したtechnical points}
\begin{itemize}[leftmargin=1.5em,itemsep=0.35em]
\item \textbf{一次情報で確認できる事実}: Metaの一次資料により、Chameleonについて2024-05-16にPaper First Submission、2024-06-18に研究公開を確認した。 対象event近傍の一次資料から image generation / 7B parameter scale / 34B parameter scale を確認できる。
\end{itemize}
\begin{minipage}{\linewidth}
% reader-facing Technical Notes generic boundary/source tail group
\begin{samepage}
{\bfseries 一次資料}
\begingroup\sloppy
\Urlmuskip=0mu plus 2mu\relax
\begin{itemize}[leftmargin=1.5em,itemsep=0.25em]
\item {\scriptsize\url{http://arxiv.org/abs/2405.09818v2}}
\item {\scriptsize\url{https://ai.meta.com/blog/meta-fair-research-new-releases/}}
\end{itemize}
\endgroup
\end{samepage}
\end{minipage}
\end{minipage}
\endgroup
\end{technicalnote}
\medskip
